\section{An endpoint criterion for local equality}

The following Hilbert-space argument explains how the off-line zero
spaces replace full half-line density in the cutoff computation.
It requires qualitative approximation at each fixed frequency window,
and makes no assertion uniform in the height of that window.

Let $\mathcal H=L^2(\mathbb R,dt)$, let
$(\tau_r f)(t)=f(t-r)$, and put
$\mathcal H_-=L^2(( -\infty,0])$ and
$\mathcal H_+=L^2([0,\infty))$.
Fix a bounded smooth compactly supported multiplier $a$, and write
$A_L=1_{[-L,L]}C_a$.  Let $Q_L$ be orthogonal projections.
For each sign choose a closed subspace
$M_+\subset\mathcal H_-$ and $M_-\subset\mathcal H_+$, and let
$K_+=\mathcal H_-\ominus M_+$ and
$K_-=\mathcal H_+\ominus M_-$.
The subscripts indicate the physical endpoint, not the sign of the
half-line.  Denote their orthogonal projections by $P_{K_\pm}$.

\begin{theorem}\label{thm:endpoint-local-criterion}
Suppose that the following conditions hold.

\begin{enumerate}
\item There are finite-rank orthogonal input projections $P_L$ such that
$\tau_{-L}P_L\tau_L$ and $\tau_LP_L\tau_{-L}$ converge strongly to zero,
and
\[
 \|(1-Q_L)A_LP_L\|_2^2\longrightarrow\sigma_a.
\]
\item There are fixed Hilbert--Schmidt operators $B_+,B_-$ whose ranges
lie respectively in $\mathcal H_-,\mathcal H_+$, operators $F_L$, and
Hilbert--Schmidt errors $r_L\to0$, such that
\begin{equation}\label{eq:endpoint-decomposition}
 A_L(1-P_L)=F_L+\tau_LB_+\tau_{-L}
                  +\tau_{-L}B_-\tau_L+r_L,
\end{equation}
with $\|(1-Q_L)F_L\|_2\to0$ and
$\|F_L\|_2\leq C(1+L)^d$ for some fixed $d$.
\item For every $v\in M_+$ and $v\in M_-$, respectively,
\[
 \|(1-Q_L)\tau_Lv\|\to0,\qquad
 \|(1-Q_L)\tau_{-L}v\|\to0.
\]
There are dense increasing unions of finite-dimensional subspaces
$K_{+,n}\subset K_+$ and $K_{-,n}\subset K_-$ for which, for every
fixed $n$ and every $N>0$,
\[
 \|Q_L\tau_LP_{K_{+,n}}\|=O_{n,N}(L^{-N}),\qquad
 \|Q_L\tau_{-L}P_{K_{-,n}}\|=O_{n,N}(L^{-N}).
\]
\item $P_{K_+}C_a$ and $P_{K_-}C_a$ are Hilbert--Schmidt.
\end{enumerate}
Then
\begin{equation}\label{eq:endpoint-local-conclusion}
 \lim_{L\to\infty}\|(1-Q_L)A_L\|_2^2
 =\sigma_a+\|P_{K_+}C_a\|_2^2+\|P_{K_-}C_a\|_2^2.
\end{equation}
\end{theorem}

\begin{proof}
Contractivity and density extend the stated convergence on finite
$K_{\pm,n}$ to every fixed vector in $K_\pm$.  Combining it with the
approximation on $M_\pm$ shows, on each corresponding half-line,
\[
 \tau_{-L}(1-Q_L)\tau_L\longrightarrow P_{K_+},\qquad
 \tau_L(1-Q_L)\tau_{-L}\longrightarrow P_{K_-}
\]
strongly.  Composition on the right by a fixed Hilbert--Schmidt
operator upgrades this to Hilbert--Schmidt convergence.  Indeed
first approximate that operator by a finite-rank operator and use
the uniform bound on the operators on the left.
Applying this observation to \eqref{eq:endpoint-decomposition} gives
\begin{equation}\label{eq:endpoint-residual}
 (1-Q_L)A_L(1-P_L)
 =\tau_LP_{K_+}B_+\tau_{-L}
  +\tau_{-L}P_{K_-}B_-\tau_L+o_{\mathcal S_2}(1).
\end{equation}

We next determine the two projected boundary operators; a bound on
their norms alone would not identify the limit.
Fix $n$ and translate \eqref{eq:endpoint-decomposition} to the upper
endpoint.  Apply $P_{K_{+,n}}$ on the left.  The bulk term tends to
zero in Hilbert--Schmidt norm, because
\[
 \begin{split}
 \|P_{K_{+,n}}\tau_{-L}F_L\tau_L\|_2
 &\leq\|Q_L\tau_LP_{K_{+,n}}\|\,\|F_L\|_2
          +\|(1-Q_L)F_L\|_2\longrightarrow0.
 \end{split}
\]
Choose $N>d$ in the finite-packet estimate for the first term.
The opposite boundary also vanishes after this finite output
projection.  To see this, choose an orthonormal basis $w_1,\ldots,w_m$
of $K_{+,n}$.  Its squared Hilbert--Schmidt norm is
\[
 \sum_{j=1}^m\|B_-^*\tau_{2L}w_j\|^2\longrightarrow0.
\]
Translations of a fixed $L^2$ vector converge weakly to zero, and
$B_-^*$ is compact.  The right input translation does not change
this norm.

The translated left side, after the finite output projection,
converges to $P_{K_{+,n}}C_a$ in Hilbert--Schmidt norm.
Indeed, put $S_L^+=1_{[-2L,0]}$ and $P_L^+=\tau_{-L}P_L\tau_L$.
The translated operator is $S_L^+C_a(1-P_L^+)$, whose adjoint is
$(1-P_L^+)C_a^*S_L^+$. For $w\in K_{+,n}\subset\mathcal H_-$,
$S_L^+w\to w$ and
\[
 P_L^+C_a^*S_L^+w
 =P_L^+C_a^*(S_L^+w-w)+P_L^+C_a^*w\longrightarrow0.
\]
The first term tends to zero by boundedness, and the second by strong
convergence on the fixed vector $C_a^*w$. Thus the adjoint tends to
$C_a^*$ on each vector of the finite output space. Summing over an
orthonormal basis gives the asserted Hilbert--Schmidt convergence.
We conclude that
\[
                 P_{K_{+,n}}B_+=P_{K_{+,n}}C_a.
\]
Letting $n\to\infty$ gives
$P_{K_+}B_+=P_{K_+}C_a$. The limit in $L$ has been taken for each
fixed packet before this exhaustion; no bound uniform in the packet
is being used. Strong convergence of the model projections
and the Hilbert--Schmidt hypothesis justify the passage in that
ideal norm as well.  The lower-endpoint argument is identical.

The two terms on the right of \eqref{eq:endpoint-residual} become
orthogonal in Hilbert--Schmidt inner product as $L\to\infty$.
For a direct proof, truncate the output of each fixed
Hilbert--Schmidt operator to a bounded interval.  Their translated
output intervals are disjoint for sufficiently large $L$.
The discarded Hilbert--Schmidt norms tend to zero uniformly in the
translations, which proves the assertion by Cauchy--Schwarz.
It follows that
\[
 \|(1-Q_L)A_L(1-P_L)\|_2^2\longrightarrow
       \|P_{K_+}C_a\|_2^2+\|P_{K_-}C_a\|_2^2.
\]
The input subspaces $P_L\mathcal H$ and $(1-P_L)\mathcal H$ are
orthogonal, so their squared Hilbert--Schmidt norms add exactly.
Adding the limit in condition 1 proves the theorem.
\end{proof}

\begin{lemma}[The critical window packet]
\label{lem:critical-window-packet}
Let $P_{0,L}$ project onto the fixed finite family of jets
\[
 v_{\gamma,j,L}(t)=(2L)^{-1/2}1_{[-L,L]}(t)e^{-i\gamma t}
              \left(\frac{t+L}{2L}\right)^j,
       \qquad 0\leq j<m_\gamma.
\]
For every $a\in C_c^\infty(\R)$,
$\|(C_a-a(\gamma))v_{\gamma,j,L}\|^2=O_{a,j}(L^{-1})$.
The projections $\tau_{\mp L}P_{0,L}\tau_{\pm L}$ tend strongly
to zero. If $\|Q_LP_{0,L}\|\to0$, then
\[
 \|(1-Q_L)S_LC_aP_{0,L}\|_2^2
                    \longrightarrow\sum_\gamma m_\gamma|a(\gamma)|^2.
\]
\end{lemma}
\begin{proof}
For $f_j(x)=1_{[0,1]}(x)x^j$, Fourier scaling gives
\[
 \widehat v_{\gamma,j,L}(\tau)
 =\sqrt{2L}\,e^{-iL(\tau-\gamma)}
                         \widehat f_j(2L(\tau-\gamma)).
\]
One integration by parts bounds $|\widehat f_j(\xi)|$ by
$C_j(1+|\xi|)^{-1}$. Boundedness and Lipschitz continuity of $a$
therefore give
\begin{align*}
 \|(C_a-a(\gamma))v_{\gamma,j,L}\|^2
 &=\int\left|a\left(\gamma+\frac\xi{2L}\right)-a(\gamma)\right|^2
                  |\widehat f_j(\xi)|^2\frac{d\xi}{2\pi}\\
 &\leq C\int\frac{\min(1,\xi^2/L^2)}{1+\xi^2}d\xi
 =O(L^{-1}).
\end{align*}
Let $V_L$ be the finite column operator of these jets and
$G_L=V_L^*V_L$. Its same-ordinate block is exactly
$(1/(j+k+1))_{j,k}$; distinct-ordinate entries are $O(L^{-1})$
by integration by parts. Thus $G_L$ tends to a positive definite
block diagonal matrix $G_\infty$ and its inverse stays bounded.
If $D_a$ is the diagonal matrix with $a(\gamma)$ repeated
$m_\gamma$ times, the preceding estimate and $S_LV_L=V_L$ give
\[
 (1-Q_L)S_LC_aV_L=V_LD_a+o_{\mathcal S_2}(1).
\]
Since $V_LG_L^{-1/2}$ is an orthonormal column, the desired squared
norm is
\[
 \Tr(G_L^{-1/2}D_a^*G_LD_aG_L^{-1/2})+o(1)
             \longrightarrow\sum_\gamma m_\gamma|a(\gamma)|^2.
\]
Here $D_a$ is scalar on each block of $G_\infty$; no exact
commutation with the finite projection was used. Finally each
translated column has amplitude $O(L^{-1/2})$ on a fixed compact set.
The bounded inverse Gram matrix shows that its projection tends to
zero on every compactly supported $L^2$ vector. Density and
contractivity extend this to strong convergence on all of $L^2$.
\end{proof}

\paragraph{Verification needed for the arithmetic application.}
For a fixed compact frequency support, the input projection $P_L$
is the span of the finitely many critical-line zero jets in a
slightly larger interval.  Its scalar contribution is
$\sigma_a=\sum_{\Re\rho=1/2}m_\rho|a(\Im\rho)|^2$.
The finite-pole inverse of $a(\tau)/L(1/2+i\tau,\chi)$, with a fixed
boundary prescription at its real poles, has a Schwartz positive
tail and a finite exponential-polynomial negative tail plus a
Schwartz remainder.  This assertion concerns only its real poles;
zeros away from the real frequency axis do not require an RH
assumption here.  The real jet constraints cancel the polynomial
tail exactly.  Smooth raw time and frequency cutoffs then give
\eqref{eq:endpoint-decomposition}; the additional term
$S_LC_a(1-S_L)$ must be included among the endpoint operators.
Omitting this term would replace $P_KC_a$ by only one input-half-line
compression and give the wrong answer.

The bulk has polynomial growth in $L$: the finite-pole inverse has
that bound in every fixed logarithmic Sobolev norm, while smooth
time and band multipliers are uniformly bounded on these norms.
The interior prolate approximation gives its Hilbert--Schmidt
distance to the finite theta image tending to zero.
The finite-packet estimates in condition 3 are precisely the
off-line estimates in the preceding section; the exponential
estimates at the upper endpoint imply every logarithmic rate.
For the lower endpoint one chooses the arbitrary logarithmic
power in the unconditional Gram estimate.

Burnol's half-line closure theorem identifies $M_+$ with a Hardy
space multiplied by the Blaschke product of the zeros to the right
of the critical line.  The lower space is obtained by the exact
theta--Fourier identity and the dual character.
His number-field proof fixes the standard finite local vectors and
varies smooth archimedean profiles with support in $(0,1)$, which
is the family required by the finite prolate approximation.
In the principal sector one must first justify the corresponding
two-constraint closure: profiles of the form
$\psi=(u\partial_u+1)\phi$, with $\phi\in C_c^\infty((0,1))$,
have zero mean and value, and Mellin multiplier
$(1-s)\zeta(s)\mathcal M\phi(s)$.
The pole at $s=1$ cancels and adds no inner zero.
The same inner-factor argument then gives the required closure,
provided its no-singular-factor and no-delay steps are carried out
for this family.

Finally, the model-space trace is the sum of its Poisson kernels,
with multiplicity.  The usual zero count makes it finite for
compact $a$.  Thus the theorem supplies local equality once these
explicit arithmetic facts have been verified.  It does not supply
a bound uniform in frequency height, and therefore alone does not
give the limit for all original compact norm tests.
Lemma~\ref{lem:local-halfline-closure} and
Theorem~\ref{thm:rational-local-limit} verify the closure and
compact-frequency approximation in the next section.
